\section{1・2 月に西日本で「雪のない平野」や瀬戸内海に落とせる場所}
\label{sec:winter-sw}

本文の 3 領域は北海道・和歌山・愛媛に限られていた。
本付録では対象を関東より西・南の全域に広げ、1・2 月に次の条件を満たす打ち上げ場所を探す。

\begin{itemize}
  \item 打ち上げ場所に雪がない。
  \item 落下点が陸の平野で、人口集中地区（DID）の外で、雪がない。
\end{itemize}

海上での回収は考えない。したがって海に落ちた回は失敗として数える。
コードは \texttt{research/winter\_sw/} にある。

\subsection{データと方法}

\begin{table}[H]
\centering
\caption{使ったデータ}
\label{tab:wsw-data}
\small
\begin{tabular}{@{}lp{0.68\linewidth}@{}}
\toprule
風 & JRA-3Q（本文と同じ製品）。北緯 24.0--37.5 度、東経 122.0--140.5 度を新たに取得した。
     2016--2025 年の 1・2 月、6 時間ごと \\
地形 & NOAA NCEI の DEM（約 \SI{500}{\metre} 間隔） \\
市街地 & 国土数値情報「人口集中地区データ」2020 年（国土交通省） \\
積雪 & ERA5 の積雪深（水当量）\citep{hersbach2020era5}。毎日 12 時（日本時間）、\ang{0.25} 格子 \\
\bottomrule
\end{tabular}
\end{table}

落下点は本文と同じ閉形式（式\eqref{eq:drift}）で、解析 1 回ごとに求めた。
1・2 月の 10 年ぶんで 2360 回の解析があり、打ち上げ場所ごとにその数だけ落下点が得られる。
落下点の分布に正規分布を仮定する必要はない。

打ち上げ場所の候補は、\ang{0.1} 間隔に並べた陸上の点 2151 か所である。
風はそれぞれの最寄りの JRA-3Q 格子点のものを使った。
あわせて、各地の代表地点 25 か所も同じ方法で評価した。

1 回の打ち上げを「成功」と数える条件は、次のすべてを満たすことである。

\begin{itemize}
  \item 落下点が\textbf{平野}：標高 \SI{100}{\metre} 以下、傾斜 \ang{2} 以下の陸地。
  \item 落下点が \textbf{DID の外}。
  \item その日、落下点と打ち上げ場所の\textbf{どちらにも雪がない}：積雪深（水当量）\SI{1}{\milli\metre} 未満。
\end{itemize}

雪の判定には、打ち上げと同じ日の積雪を使った。
10 年ぶんの平均ではなく、その日ごとに判定している。

\begin{figure}[H]
\centering
\includegraphics[width=\linewidth]{winter_sw/W1_target_and_snow.pdf}
\caption[落下の目標と雪のない日の割合]{左は落下の目標（緑）と DID（赤）。
陸地のうち平野は 20 \% で、平野の 12 \% が DID に入る。
右は 1・2 月に雪のない日の割合。濃いほど雪の日が多い。}
\label{fig:wsw-base}
\end{figure}

\subsection{1・2 月の落下位置の傾向}

\begin{figure}[H]
\centering
\includegraphics[width=\linewidth]{winter_sw/W2_drift_jan_feb.pdf}
\caption[1・2 月の平均落下位置]{標準の飛行条件（上昇 \SI{5}{\metre\per\second}、
破裂 \SI{30}{\kilo\metre}、降下 \SI{6}{\metre\per\second}）での平均落下位置。
点が打ち上げ地点、矢先が平均落下点で、矢印は実寸である。}
\label{fig:wsw-drift}
\end{figure}

1・2 月は、どこから上げても\textbf{ほぼ真東}（方位 \ang{82}--\ang{89}）へ流される。
飛行距離は本土で \SIrange{240}{270}{\kilo\metre}、南ほど短く、
奄美で \SI{233}{\kilo\metre}、那覇で \SI{197}{\kilo\metre}、石垣で \SI{163}{\kilo\metre} である。
冬の亜熱帯ジェットがこの緯度帯の真上にあるためで、
第 \ref{sec:wakayama} 章で見た和歌山の冬と同じ性質が、西日本全体に広がっている。

1 月前半・1 月後半・2 月前半・2 月後半の違いは小さい。
本土では 2 月前半が最も遠く（例：大阪 \SI{271}{\kilo\metre}）、
2 月後半が最も近い（同 \SI{247}{\kilo\metre}）が、差は約 \SI{25}{\kilo\metre} で、向きはほとんど変わらない。
荒れた日（p90）には、本土で \SIrange{315}{345}{\kilo\metre} に達する。

\subsection{標準の飛行条件での結果}

\begin{figure}[H]
\centering
\includegraphics[width=\linewidth]{winter_sw/W3_suitability_standard.pdf}
\caption[標準の飛行条件で目標に落ちる割合]{打ち上げ場所ごとに、
「雪がなく、DID の外の平野」に落ちる割合を半月ごとに示したもの。
打ち上げ場所に雪がある日と、海に落ちた日は失敗に数えている。}
\label{fig:wsw-std}
\end{figure}

\begin{table}[H]
\centering
\caption{代表地点から標準の飛行条件で上げた場合（1・2 月の平均）。
「目標」は雪がなく DID の外の平野に落ちる割合、「平野（DID 含む）」は DID の中も平野として数えた場合の値、「海」は海に落ちる割合で、
「打ち上げ地に雪」は打ち上げ場所に雪がある日の割合である。}
\label{tab:wsw-sites}
\small
\setlength{\tabcolsep}{4pt}
\begin{tabular}{@{}p{0.3\linewidth}rrrrr@{}}
\toprule
地点 & {距離 [\si{\kilo\metre}]} & {目標} & {平野（DID 含む）} & {海} & {打ち上げ地に雪} \\
\midrule
岡山 & 260 & \textbf{13 \%} & 17 \% & 32 \% & 14 \% \\
名古屋 & 261 & \textbf{12 \%} & 21 \% & 52 \% & 13 \% \\
松江 & 253 & 6 \% & 7 \% & 27 \% & 17 \% \\
大阪 & 262 & 5 \% & 7 \% & 57 \% & 18 \% \\
広島 & 259 & 5 \% & 9 \% & 31 \% & 36 \% \\
鳥取 & 253 & 5 \% & 7 \% & 7 \% & 50 \% \\
徳島・和歌山・津 & 263--265 & 4 \% & 5--7 \% & 73--85 \% & 10--11 \% \\
福岡・松山 & 259--262 & 3 \% & 4--5 \% & 35--44 \% & 7--18 \% \\
浜松 & 266 & 3 \% & 3 \% & 92 \% & 6 \% \\
高知・大分・熊本 & 261--265 & 1 \% & 1--2 \% & 77--82 \% & 2--5 \% \\
愛南・室戸・潮岬・宮崎・鹿児島 & 263--269 & 0 \% & 0 \% & 94--100 \% & 0--3 \% \\
種子島・奄美・那覇・石垣 & 163--257 & 0--1 \% & 0--1 \% & 98--100 \% & 0 \% \\
金沢 & 246 & 0 \% & 0 \% & 3 \% & 97 \% \\
\bottomrule
\end{tabular}
\end{table}

結果ははっきりしている。
\textbf{標準の飛行条件では、1・2 月を通した平均で目標に落ちる割合が 25 \% を超える場所は無い}。
平均で最も高いのは岐阜県南部（北緯 35.6 度、東経 137.1 度付近）の 22 \%（半月ごとに 18--24 \%）で、
落下先は関東平野である。
半月ごとに見ると、2 月後半の長野県（北緯 36.0 度、東経 137.9 度付近）の 30 \% が最大である。
代表地点では岡山（13 \%）と名古屋（12 \%）が最も高い。

理由は 2 つある。

\paragraph{東へ 250 km 先に、大きな平野が少ない。}
九州・四国の南岸・紀伊半島・東海から上げると、落下点のほとんどが太平洋に出る
（愛南・潮岬・宮崎などは 94--100 \%）。
中国地方や近畿の内陸から上げると、今度は中部・近畿の山地に落ちる。
250 km 東に平野があるのは、岐阜・名古屋（→ 関東平野）、岡山（→ 濃尾平野・伊勢湾岸）などに限られる。

\paragraph{その平野は大都市圏である。}
250 km 東にある平野の多くは、関東・濃尾・大阪の大都市圏で、DID が広い。
DID を除くと、名古屋では 21 \% が 12 \% に、岐阜県南部では 30--34 \% が 18--24 \% に下がる。

日本海側は、打ち上げ場所の雪で失敗する日が多い。鳥取で約 5 割、金沢で 97 \% である。

\subsection{飛行条件を変えた場合}

飛行距離が短くなれば、落下点は打ち上げ場所の近くの平野に入りやすくなる。
そこで上昇速度と破裂高度を変えた 5 つの条件で、同じ評価を行った。
降下速度は海面相当 \SI{6}{\metre\per\second} で共通にした。

\begin{figure}[H]
\centering
\includegraphics[width=\linewidth]{winter_sw/W4_scenarios.pdf}
\caption[飛行条件ごとの比較]{飛行条件ごとの「目標に落ちる割合」（1・2 月の平均）。
丸は最も割合が高い地点。見出しの飛行距離は、打ち上げ候補全体の平均である。}
\label{fig:wsw-scen}
\end{figure}

\begin{table}[H]
\centering
\caption{飛行条件ごとの飛行距離と、最良地点で目標に落ちる割合（1・2 月の平均）}
\label{tab:wsw-scen}
\begin{tabular}{@{}lrrr@{}}
\toprule
飛行条件 & {飛行距離 [\si{\kilo\metre}]} & {最良地点} & {目標に落ちる割合} \\
\midrule
標準（上昇 5 m/s・破裂 30 km） & 255 & 岐阜県南部 & 22 \% \\
低い破裂（5 m/s・20 km） & 239 & 長野県 & 23 \% \\
速い上昇（8 m/s・30 km） & 188 & 長野県 & 36 \% \\
速い上昇・破裂 25 km & 187 & 長野県 & 37 \% \\
\textbf{速い上昇・破裂 20 km} & \textbf{177} & \textbf{長野県} & \textbf{41 \%} \\
\bottomrule
\end{tabular}
\end{table}

破裂高度を下げるだけでは、飛行距離はあまり縮まない（\SI{255}{\kilo\metre} → \SI{239}{\kilo\metre}）。
第 \ref{sec:wakayama} 章の最後で見たとおり、冬の落下距離の約 6 割は
高度 \SIrange{8}{16}{\kilo\metre} のジェットの層で生じるので、
その層を速く通り抜けるほうが効果が大きい。
上昇速度を \SI{8}{\metre\per\second} にすると、飛行距離は約 7 割（\SI{188}{\kilo\metre}）になる。

\begin{figure}[H]
\centering
\includegraphics[width=\linewidth]{winter_sw/W5_suitability_best.pdf}
\caption[速い上昇・破裂 20 km で目標に落ちる割合]{最も成績のよい飛行条件
（上昇 \SI{8}{\metre\per\second}、破裂 \SI{20}{\kilo\metre}）での、半月ごとの結果。}
\label{fig:wsw-best}
\end{figure}

最も成績のよい条件（上昇 \SI{8}{\metre\per\second}、破裂 \SI{20}{\kilo\metre}）では、
長野県の佐久〜八ヶ岳付近（北緯 36.0--36.1 度、東経 138.2--138.4 度）から上げた場合に、
37--45 \% が目標に落ちる。
落下先は北関東の平野（埼玉・群馬・茨城の DID の外）である。

\begin{figure}[H]
\centering
\includegraphics[width=\linewidth]{winter_sw/W6_landing_clouds.pdf}
\caption[落下点の分布]{最良地点から上げた場合の、解析 1 回ごとの落下点。
上は標準の飛行条件、下は上昇 \SI{8}{\metre\per\second}・破裂 \SI{20}{\kilo\metre}。
青が目標に落ちた回、灰色がそれ以外（DID・山地・海、または雪の日）。}
\label{fig:wsw-clouds}
\end{figure}

\begin{table}[H]
\centering
\caption{地方ごとの最良地点（1・2 月の平均）。地方は、海岸線で本州・九州・四国と島を分けたうえで、
本州の中を経度・緯度で分けた。「最良の条件」は 5 つの飛行条件のうち最もよいもの。}
\label{tab:wsw-areas}
\small
\begin{tabular}{@{}lp{0.42\linewidth}rr@{}}
\toprule
地方 & 最良地点（標準条件 / 最良の条件） & {標準条件} & {最良の条件} \\
\midrule
中部（長野・岐阜・山梨・北陸） & 岐阜県南部 / 長野県 佐久〜八ヶ岳 & 22 \% & \textbf{41 \%} \\
東海（静岡・愛知・三重） & 愛知県北部 / 山梨・静岡県境（富士川沿い） & 17 \% & 20 \% \\
近畿 & 兵庫県南西部 / 兵庫県中部 & 12 \% & 19 \% \\
中国 & 岡山県 & 14 \% & 15 \% \\
四国 & 香川県 & 9 \% & 9 \% \\
九州 & 福岡県北九州 / 長崎県北部 & 7 \% & 7 \% \\
離島 & 小豆島 / 五島列島 & 12 \% & 15 \% \\
南西諸島 & 先島・沖縄本島 & 3 \% & 4 \% \\
\bottomrule
\end{tabular}
\end{table}

中部以外の地方では、飛行条件を変えても 20 \% 前後かそれ以下にとどまる。
四国と九州の本土はどの条件でも 10 \% に届かない。
東へ 180--260 km の範囲が、ほとんど海（太平洋や瀬戸内海）か山地になるためである。
離島では、上昇を速くすると五島列島から佐賀・熊本の平野に落とせる日が増え、15 \% になる。
南西諸島は、どの条件でも落下点のほとんどが海になる。

\subsection{瀬戸内海への着水を認める場合}
\label{sec:wsw-seto}

ここまでは海に落ちた回をすべて失敗に数えた。
波の穏やかな瀬戸内海であれば船で回収できるので、瀬戸内海への着水も成功に数えて評価し直した。
瀬戸内海の範囲は、関門海峡・豊予海峡（関崎〜佐田岬）・鳴門海峡・紀淡海峡を境とする内側の海とし、
外洋に近い紀伊水道と豊後水道は含めない（別府湾は含める）。
着水した回は落下点の積雪を問わないが、打ち上げ場所に雪がある日は引き続き失敗に数える。
コードでは \texttt{python winter\_sw/analyze.py --allow-seto} で切り替わる。

\begin{figure}[H]
\centering
\includegraphics[width=\linewidth]{winter_sw_seto/W3_suitability_standard.pdf}
\caption[瀬戸内海を含めた場合の割合（標準の飛行条件）]{瀬戸内海への着水も成功に数えた場合の、
標準の飛行条件での「目標に落ちる割合」。色の目盛りは図\ref{fig:wsw-std} と違い、70 \% までである。}
\label{fig:wsw-seto-std}
\end{figure}

結果は大きく変わる。
1・2 月は真東へ約 \SI{250}{\kilo\metre} 流されるので、
九州北西部や山口県西部から上げると、ちょうど周防灘から伊予灘・広島湾に落ちる。

\begin{table}[H]
\centering
\caption{瀬戸内海への着水を認めた場合の、地方ごとの最良地点（1・2 月の平均）}
\label{tab:wsw-seto-areas}
\small
\begin{tabular}{@{}lp{0.42\linewidth}rr@{}}
\toprule
地方 & 最良地点（標準条件 / 最良の条件） & {標準条件} & {最良の条件} \\
\midrule
離島 & 壱岐 & \textbf{47 \%} & \textbf{60 \%} \\
九州 & 佐賀県北部（唐津・呼子付近） & 37 \% & \textbf{62 \%} \\
中国 & 山口県西部（長門付近） & 39 \% & 47 \% \\
中部（長野・岐阜・山梨・北陸） & 岐阜県南部 / 長野県 佐久〜八ヶ岳 & 22 \% & 41 \% \\
東海（静岡・愛知・三重） & 愛知県北部 / 山梨・静岡県境 & 17 \% & 20 \% \\
四国 & 愛媛県北部（松山の北） & 10 \% & 19 \% \\
近畿 & 兵庫県南西部 / 兵庫県中部 & 12 \% & 19 \% \\
南西諸島 & 先島 & 3 \% & 4 \% \\
\bottomrule
\end{tabular}
\end{table}

\paragraph{壱岐は標準の飛行条件のままで 47 \% に達する。}
半月ごとに 40--60 \% で、2 月後半が最も高い。
成功した回のほとんど（全体の 36--55 \%）が瀬戸内海への着水で、陸の平野に落ちるのは 5 \% 前後である。
1・2 月に打ち上げ場所に雪がある日は、10 年を通して無かった。
ただし壱岐のような小さな島では、ERA5 の格子のほとんどが海になるので、
積雪が少なめに出ている可能性がある（\ref{sec:wsw-notes} 節）。

\paragraph{上昇を速くすると、九州北部の本土から 6 割に届く。}
上昇 \SI{8}{\metre\per\second} にすると飛行距離が約 \SI{190}{\kilo\metre} に縮み、
佐賀県北部（唐津・呼子付近）から上げた場合に 55--66 \% が目標に入る。
打ち上げ場所に雪がある日は 1--5 \% である。

\begin{figure}[H]
\centering
\includegraphics[width=\linewidth]{winter_sw_seto/W6_landing_clouds.pdf}
\caption[瀬戸内海を含めた場合の落下点]{上は壱岐から標準の飛行条件で、
下は唐津付近から上昇 \SI{8}{\metre\per\second}・破裂 \SI{20}{\kilo\metre} で上げた場合の、解析 1 回ごとの落下点。
水色が瀬戸内海、青い点が目標に入った回である。
外れた回は、手前の北九州・山口の山地か、瀬戸内海を越えた四国の山地、または豊後水道側の海に落ちている。}
\label{fig:wsw-seto-clouds}
\end{figure}

瀬戸内海を目標に含めても、打ち上げ場所の選択肢は西日本の西の端に限られる。
東へ \SI{250}{\kilo\metre} 先が瀬戸内海になるのは、九州北西部・壱岐・山口県西部だけである。
それより東から上げると、瀬戸内海を通り越して近畿・中部の山地や太平洋に落ちる。

\subsection{読むときの注意}
\label{sec:wsw-notes}

\begin{itemize}
  \item \textbf{瀬戸内海は船の往来が多い。}
        着水を認める場合は、回収船の手配に加えて、航路との位置関係を事前に確認する必要がある。
        また瀬戸内海の境界は本付録で決めたもので、法令上の「瀬戸内海」より狭い。
  \item \textbf{「DID の外の平野」は、回収しやすい場所と同じではない。}
        北関東の平野は DID の外でも人家や道路が多い。
        また首都圏の上空は、成田・羽田の空港に近い空域にかかる。
        打ち上げの届出や飛行の可否は、この統計とは別に確認する必要がある。
  \item \textbf{ERA5 の積雪は \ang{0.25}（約 \SI{25}{\kilo\metre}）格子である。}
        沿岸の都市でも、同じ格子に入る山地の雪を拾う。
        広島（36 \%）や大阪（18 \%）の「打ち上げ地に雪」は、実際より多めに出ている可能性が高い。
        逆に壱岐のような小さな島では、格子のほとんどが海になるので少なめに出る。
  \item \textbf{平野の判定は標高と傾斜だけで行っている。}
        湖（琵琶湖など、湖面の標高が \SI{1}{\metre} を超えるもの）、河川、森林、農地を区別していない。
  \item \textbf{最良地点の長野県は標高 \SIrange{700}{1300}{\metre} の高地である。}
        計算では地上高度を 0 としているので、実際の飛行時間はわずかに短くなる。
        また 1・2 月に打ち上げ場所に雪がある日が 12--37 \% ある。
  \item \textbf{風は打ち上げ時刻・地点で固定している。}
        本文の第 \ref{sec:fidelity} 章のとおり、この近似による誤差は冬で \SIrange{5}{6}{\kilo\metre} で、
        結論は変わらない。ただし当日の予報誤差は含まれていない。
\end{itemize}

\subsection{まとめ}

\begin{itemize}
  \item 1・2 月は西日本のどこから上げても、真東へ \SIrange{240}{270}{\kilo\metre}（南西諸島は \SIrange{160}{230}{\kilo\metre}）流される。
        前半と後半の違いは小さい。
  \item 陸の平野（DID の外、雪なし）だけを目標にすると、標準の飛行条件では 1・2 月の平均でどこでも 25 \% 未満である。
        上昇を速くして飛行距離を 7 割に縮めても、長野県の佐久〜八ヶ岳付近から北関東の平野へ落とす形で約 4 割が上限になる。
  \item 瀬戸内海への着水を認めると、九州北西部と壱岐が有力になる。
        壱岐は標準の飛行条件のままで 47 \%、佐賀県北部は上昇を速くすると 62 \% が目標に入る。
        どちらも 1・2 月に打ち上げ場所の雪がほとんど無い。
\end{itemize}

地形は NOAA NCEI の DEM、市街地は国土数値情報（人口集中地区データ、国土交通省）、
積雪は ERA5（Copernicus Climate Change Service の情報を加工したものを含む）による。
